\documentclass[11pt,a4paper]{article}
% Préambule commun aux documents de PythonIPM/documentation.
% Compilation : tectonic <fichier>.tex
% tectonic compile avec XeTeX, qui lit l'UTF-8 nativement : surtout PAS inputenc ni fontenc,
% faits pour pdfLaTeX. Ce sont eux qui décodaient mal les caractères multi-octets — « cœur »
% ressortait en « cur », et les tirets cadratins en caractères parasites.
\usepackage{lmodern}
\usepackage[french]{babel}
\usepackage{microtype}
\usepackage[a4paper,margin=2.4cm]{geometry}
\usepackage{amsmath,amssymb}
\usepackage{booktabs}
\usepackage{longtable}
\usepackage{array}
\usepackage{ragged2e}
\usepackage{xcolor}
\usepackage{tcolorbox}
\usepackage{enumitem}
\usepackage{fancyhdr}
\usepackage[hidelinks]{hyperref}

\definecolor{encadre}{HTML}{F4F6F8}
\definecolor{bordure}{HTML}{4A6FA5}
\definecolor{alerte}{HTML}{B03A2E}

% Encadré « à retenir » : la lecture non technique du passage qui précède
\newtcolorbox{retenir}[1][]{
  colback=encadre, colframe=bordure, boxrule=0.6pt, arc=2pt,
  left=8pt, right=8pt, top=6pt, bottom=6pt,
  fonttitle=\bfseries, title={#1}
}

% Encadré d'avertissement : un point ouvert, un écart assumé
\newtcolorbox{attention}[1][]{
  colback=white, colframe=alerte, boxrule=0.6pt, arc=2pt,
  left=8pt, right=8pt, top=6pt, bottom=6pt,
  % le bandeau de titre a déjà la couleur d'alerte en fond : un titre `coltitle=alerte`
  % écrirait du rouge sur du rouge. On garde le blanc par défaut, comme l'encadré « retenir ».
  fonttitle=\bfseries, title={#1}
}

\newcommand{\col}[1]{\texttt{#1}}
% \ensuremath et non \textsubscript : \ipm est employé aussi bien dans le texte que dans les
% formules, et « M\textsubscript{0} » placé en mode mathématique est une erreur — silencieuse en
% PDF, mais qui fait échouer la conversion vers Word.
\newcommand{\ipm}{\ensuremath{M_{0}}}

\setlist[itemize]{itemsep=2pt, topsep=4pt}
\setlist[description]{style=nextline, leftmargin=1.2cm, itemsep=5pt}

% La source du document apparaît en tête de chaque page. Les documents RGPH la redéfinissent
% par \renewcommand{\source}{...} AVANT \begin{document}.
\newcommand{\source}{EHCVM 2021}

\pagestyle{fancy}
\fancyhf{}
\fancyhead[L]{\small\itshape IPM Côte d'Ivoire --- \source}
\fancyhead[R]{\small\thepage}
\renewcommand{\headrulewidth}{0.3pt}

\renewcommand{\arraystretch}{1.15}


\title{\bfseries Note sur la dimension Santé\\[4pt]
\large Pourquoi l'assurance maladie ne peut pas la porter seule}
\author{Pipeline \col{PythonIPM} --- IPM Côte d'Ivoire, EHCVM 2021}
\date{Août 2026}

\begin{document}
\maketitle

\begin{abstract}
\noindent La dimension Santé de l'IPM ivoirien a d'abord reposé sur un unique indicateur ---
l'absence de couverture maladie --- faute de pouvoir mesurer la mortalité juvénile dans l'EHCVM.
Cette note montre pourquoi ce montage produisait un indice non plausible, expose le mécanisme
arithmétique en cause, documente la décision prise (l'ajout de l'insécurité alimentaire mesurée
par l'échelle FIES de la FAO) et en mesure l'effet. Elle se termine sur ce qui reste ouvert.
\end{abstract}

\section{Le point de départ}

Le tableau de référence national place dans la dimension Santé un indicateur de \textbf{mortalité
juvénile}. L'EHCVM 2021 ne permet pas de le construire : l'enquête n'a ni module décès ni
historique des naissances. La base des chocs enregistre bien un « décès d'un membre du ménage »
--- 1\,328 ménages sur trois ans, avec la date --- mais \textbf{sans l'âge du défunt}. Il est donc
impossible d'isoler les décès d'enfants de moins de dix-huit ans.

La dimension a d'abord été mesurée par un seul indicateur de substitution : \emph{aucun membre du
ménage n'est couvert par une assurance maladie}. Le choix se défendait --- la question est
renseignée pour 100~\% des ménages, sans biais de sélection, et se lit sans ambiguïté.

\section{Le problème}

Avec quatre dimensions équipondérées à 0,25, un indicateur seul dans sa dimension \textbf{porte
tout son poids}. Or 93,4~\% des ménages ivoiriens n'ont aucun membre assuré. Trois conséquences
en découlent, toutes arithmétiques et non statistiques.

\subsection{L'indicateur dominait le score}

L'assurance maladie apportait en moyenne \textbf{0,2336} sur un score moyen de 0,4603, soit
\textbf{50,8~\% du score de privation} à elle seule. Les douze autres indicateurs se partageaient
l'autre moitié.

\subsection{Elle fonctionnait comme un laissez-passer vers le seuil}

Un ménage non assuré partait déjà à 0,25 sur les 0,3333 requis pour être déclaré pauvre. Deux
privations de conditions de vie ($2 \times 0{,}0417 = 0{,}0833$) suffisaient alors à franchir
exactement le seuil. \textbf{403 ménages} se trouvaient dans ce cas précis, et 396 avaient un
score exactement égal à $k$ : un ménage sur trente basculait sur une différence de troisième
décimale.

\subsection{Le résultat n'était pas plausible}

Une simulation retirant l'indicateur et renormalisant les poids donne :

\begin{center}
\begin{tabular}{@{}lr@{}}
\toprule
\textbf{Configuration} & \textbf{$H$ (population pauvre)} \\
\midrule
Treize indicateurs, assurance seule en Santé & 79,7~\% \\
Les mêmes, sans l'assurance maladie & 45,7~\% \\
\bottomrule
\end{tabular}
\end{center}

Le second ordre de grandeur est celui publié pour la Côte d'Ivoire dans la table globale de
l'OPHI. Le premier ne l'est pas.

\begin{retenir}[Le fond du problème]
Un indicateur qui prive 93~\% de la population ne discrimine plus : il déplace le niveau de
l'indice sans rien apprendre sur \emph{qui} est pauvre. Ce n'était pas une erreur de calcul --- le
code appliquait fidèlement les choix actés --- mais la rencontre de deux décisions qui se
combinent mal : substituer l'assurance maladie à la mortalité juvénile, \emph{et} conserver quatre
dimensions équipondérées.
\end{retenir}

\section{La décision}

Un \textbf{second indicateur} a été ajouté à la dimension Santé : l'insécurité alimentaire mesurée
par l'échelle \textbf{FIES} (\emph{Food Insecurity Experience Scale}) de la FAO, présente dans la
section 8A de l'EHCVM. Huit questions oui/non portant sur les douze derniers mois, ordonnées par
sévérité croissante, dont le nombre de « oui » forme un score de 0 à 8.

\begin{center}
\small
\begin{tabular}{@{}rlp{2.6cm}@{}}
\toprule
\textbf{\#} & \textbf{Question} & \textbf{Sévérité} \\
\midrule
1 & inquiétude de manquer de nourriture & légère \\
2 & n'a pas pu manger sain et nutritif & légère \\
3 & a mangé peu varié & légère \\
4 & a sauté un repas & modérée \\
5 & a mangé moins que nécessaire & modérée \\
6 & plus de nourriture dans le ménage & modérée \\
7 & a eu faim sans manger & sévère \\
8 & est resté une journée entière sans manger & sévère \\
\bottomrule
\end{tabular}
\end{center}

Le seuil retenu est \textbf{score $\ge 4$}, qui correspond à l'insécurité alimentaire modérée ou
sévère --- le seuil standardisé de l'\textbf{indicateur ODD 2.1.2}. Il n'a donc pas fallu inventer
un seuil : celui-ci est comparable aux publications de la FAO.

\subsection{Pourquoi cet indicateur plutôt qu'un autre}

La justification est de fond, et pas seulement d'opportunité. Dans l'IPM mondial, la dimension
Santé est composée de la \textbf{nutrition} et de la mortalité juvénile. Remplacer la seconde
--- impossible à mesurer ici --- par une mesure de la première ne dénature pas la dimension :
c'est son autre pilier officiel.

S'y ajoute une raison statistique. La distribution du score FIES le sépare nettement de
l'assurance maladie :

\begin{center}
\small
\begin{tabular}{@{}lrrrrrrrrr@{}}
\toprule
\textbf{Score FIES} & 0 & 1 & 2 & 3 & 4 & 5 & 6 & 7 & 8 \\
\midrule
\textbf{Part des ménages} & 30,5 & 7,5 & 7,8 & 10,1 & 9,2 & 9,8 & 7,9 & 8,5 & 7,4 \\
\bottomrule
\end{tabular}
\end{center}

Un tiers des ménages à zéro, le reste étalé sur toute l'échelle. L'assurance maladie, elle, oppose
93,4~\% à 6,6~\% : elle sépare mal parce qu'elle ne varie presque pas. C'est exactement ce qu'on
demande à un indicateur de privation --- séparer --- et c'est ce que le FIES fait.

Le taux de privation retenu s'établit à \textbf{42,7~\% des ménages}, soit 40,8~\% de la
population.

\section{L'effet mesuré}

La dimension Santé compte désormais deux indicateurs à \textbf{0,125} chacun : le poids de
l'assurance maladie est mécaniquement divisé par deux.

\begin{center}
\begin{tabular}{@{}lrr@{}}
\toprule
& \textbf{Avant} & \textbf{Après} \\
& \small(assurance seule) & \small(assurance + FIES) \\
\midrule
Poids de l'assurance maladie & 0,2500 & 0,1250 \\
Part de l'assurance dans le score moyen & 50,8~\% & 29,4~\% \\
Score moyen $c_i$ & 0,4603 & 0,3968 \\
\textbf{$H$ --- incidence} & \textbf{79,7~\%} & \textbf{65,8~\%} \\
Vulnérables & 12,7~\% & 19,1~\% \\
Pauvreté sévère & 44,0~\% & 26,8~\% \\
Ménages au score exactement égal à $k$ & 396 & 928 \\
\bottomrule
\end{tabular}
\end{center}

L'effet est net sur la pauvreté sévère, qui passe de 44,0~\% à 26,8~\% : c'est là que la
domination d'un indicateur unique se voyait le plus.

\section{Limites assumées}

\begin{description}
  \item[Le FIES n'est pas une mesure d'état nutritionnel.] Il mesure l'insécurité alimentaire
  \emph{vécue}. L'OPHI, lui, mesure la nutrition par l'anthropométrie --- indice de masse
  corporelle des adultes, retard de croissance des enfants. L'EHCVM n'a pas de module
  anthropométrique : le FIES est le meilleur substitut disponible, pas un équivalent.

  \item[C'est du déclaratif rétrospectif sur douze mois.] Il est donc sensible à la saison de
  collecte. Les deux vagues de l'EHCVM ayant été collectées à des saisons différentes, une
  vérification de la comparabilité entre vagues reste à mener.

  \item[195 ménages (1,5~\%) n'ont pas de score complet], pour cause de non-réponse ou de refus
  sur au moins un des huit items. Ils sont traités comme non privés, conformément à la convention
  générale du pipeline. C'est un choix --- celui de ne pas pénaliser un refus de répondre --- et
  non une déduction.
\end{description}

\section{Ce qui reste ouvert}

\begin{attention}[$H = 65{,}8~\%$ demeure au-dessus de l'ordre de grandeur de référence]
L'assurance maladie pèse encore \textbf{26,0~\% de l'IPM} --- moins qu'avant, mais beaucoup pour
un indicateur qui prive 93~\% des ménages. Trois leviers restent disponibles, non tranchés à ce
jour.
\end{attention}

\begin{enumerate}
  \item \textbf{Abandonner l'assurance maladie} et ne garder que le FIES en Santé. Cohérent avec
  le constat qu'un indicateur à 93~\% n'informe pas, mais laisse la dimension sur une seule jambe
  --- avec le risque symétrique déjà observé pour l'Emploi.

  \item \textbf{Rééquilibrer les dimensions.} Rien n'oblige à quatre dimensions égales. L'IPM
  mondial en retient trois --- Santé, Éducation, Niveau de vie --- à un tiers chacune.

  \item \textbf{Revoir d'autres seuils dans le même esprit.} L'année de scolarité (72,5~\%),
  l'énergie de cuisson (76,8~\%) et les toilettes (77,3~\%) privent aussi une large majorité des
  ménages, et posent la même question de pouvoir discriminant.
\end{enumerate}

Ces arbitrages relèvent du cadrage et non de la technique : chacun se modifie en une ligne dans
\col{pipeline/02\_construction\_matrice\_situationnelle.py}.

\vfill
\noindent\rule{\linewidth}{0.4pt}\\[2pt]
\small
Documents associés : \emph{Méthodologie de calcul} et \emph{Dictionnaire des sorties}.

\end{document}
